% Generated from locale/tr/content/computability/computability-theory/ce-closed-cup-cap.tex; do not hand-edit.


\olfileid[tr]{cmp}{thy}{clo}

\olsection[Hesaplanabilir Biçimde Numaralandırılabilir Kümelerin Birleşimi ve
Kesişimi]{Hesaplanabilir Biçimde Numaralandırılabilir Kümeler Birleşim ve
Kesişim Altında Kapalıdır}

Aşağıdaki teorem hesaplanabilir biçimde numaralandırılabilir kümeler için
bazı kapanış özellikleri verir.

\begin{thm}
$A$ ile $B$ hesaplanabilir biçimde numaralandırılabilir olsun. O zaman
$A\cap B$ ile $A\cup B$ de öyledir.
\end{thm}

\begin{proof}
\olref[eqc]{thm:ce-equiv}, hesaplanabilir biçimde numaralandırılabilir
kümelerin çeşitli karakterizasyonlarını kullanmamızı sağlar. Açıklayıcı
olması için birkaç farklı ispat vereceğiz.

İlk ispatta $A$, hesaplanabilir $f$ fonksiyonuyla; $B$ ise hesaplanabilir
$g$ fonksiyonuyla numaralandırılsın. Şöyle tanımlayalım:
\begin{align*}
  h(x) & = \umin{y}{(f(y)=x\lor g(y)=x)} \text{ ve}\\
  j(x) & = \umin{y}{(f((y)_0)=x\land g((y)_1)=x)}.
\end{align*}
O zaman $A\cup B$, $h$'nin; $A\cap B$ ise $j$'nin tanım kümesidir.

\begin{explain}
Hesaplama açısından olan şudur: $A$ ile $B$'yi numaralandıran yordamlar
verildiğinde bir $x$ elemanının $A\cup B$ içinde olup olmadığına iki
numaralandırmadan birinde $x$'i arayarak yarı karar verebiliriz. $x$'in
$A\cap B$ içinde olup olmadığına ise aynı anda iki numaralandırmada da
$x$'i arayarak yarı karar verebiliriz.
\end{explain}

İkinci ispatta yine $A$, $f$ ile ve $B$, $g$ ile numaralandırılsın. Şöyle
tanımlayalım:
\[
  k(x)=\begin{cases}
    f(x/2) & \text{$x$ çiftse}\\
    g((x-1)/2) & \text{$x$ tekse.}
  \end{cases}
\]
O zaman $k$, $A\cup B$'yi numaralandırır; çünkü $f$ ile $g$'nin sunduğu
numaralandırmalar arasında sırayla geçiş yapar. $A\cap B$'yi numaralandırmak
daha güçtür. $A\cap B$ boşsa kendiliğinden hesaplanabilir biçimde
numaralandırılabilir. Değilse $c$, $A\cap B$'nin herhangi bir elemanı olsun
ve $l$'yi
\[
  l(x)=\begin{cases}
    f((x)_0) & \text{$f((x)_0)=g((x)_1)$ ise}\\
    c & \text{aksi durumda}
  \end{cases}
\]
ile tanımlayalım. Hesaplama açısından $l$, $f$ ile $g$ numaralandırmalarındaki
eleman çiftleri arasında dolaşır ve bulduğu her eşleşmeyi çıktı olarak verir;
aksi durumda $c$'yi döndürerek yerinde sayar.

Son ispatta $A$, kısmi $m(x)$ fonksiyonunun; $B$ ise kısmi $n(x)$
fonksiyonunun tanım kümesi olsun. O zaman $A\cap B$, kısmi
$m(x)+n(x)$ fonksiyonunun tanım kümesidir.

\begin{explain}
Hesaplama açısından $A$, $m$ yordamının durduğu ve $B$, $n$ yordamının
durduğu değerler kümesiyse $A\cap B$, iki yordamın da durduğu değerler
kümesidir.
\end{explain}

$m$ ile $n$'yi paralel olarak benzetmek gerektiğinden $A\cup B$'yi durma
değerleri kümesi olarak ifade etmek daha zordur. $d$, $m$'nin ve $e$,
$n$'nin bir indisi olsun; yani $m=\cfind{d}$ ve $n=\cfind{e}$. O zaman
$A\cup B$,
\[
  p(x)=\umin{y}{(T(d,x,y)\lor T(e,x,y))}
\]
fonksiyonunun tanım kümesidir.
\begin{explain}
Hesaplama açısından $p$, $x$ girdisinde $m$ ya da $n$ için duran bir hesap
arar ve ikisinden birini bulursa durur.
\end{explain}
\end{proof}

